\documentclass[11pt,a4paper]{article}
\usepackage[T1]{fontenc}
\usepackage{lmodern}
\usepackage[margin=24mm,headheight=15pt]{geometry}
\usepackage{amsmath,amssymb,amsthm,mathtools}
\usepackage{booktabs,array,tabularx}
\usepackage{microtype,fancyhdr,xcolor}
\usepackage[unicode]{hyperref}
\definecolor{ink}{RGB}{28,61,79}
\hypersetup{colorlinks=true,linkcolor=ink,citecolor=ink,urlcolor=ink,
pdftitle={SRG Geometry: Local Response, Neighbor Exchange, and Compression},
pdfauthor={Hilmir Frímann Halldórsson},
pdfsubject={v0.1 manuscript for GPT review; one-pass exact-real mathematical analysis},
pdfkeywords={SRG, geometry, directed neighbors, compression, local response},
pdfcreator={LaTeX}}
\pagestyle{fancy}\fancyhf{}
\fancyhead[L]{\small SRG Geometry}
\fancyhead[R]{\small v0.1 \enspace|\enspace Manuscript for review}
\fancyfoot[C]{\thepage}
\renewcommand{\headrulewidth}{0.3pt}
\setlength{\parindent}{0pt}
\setlength{\parskip}{5pt plus 1pt minus 1pt}
\setlength{\emergencystretch}{2em}
\clubpenalty=10000 \widowpenalty=10000
\allowdisplaybreaks[1]
\newcolumntype{Y}{>{\raggedright\arraybackslash}X}
\newtheorem{theorem}{Theorem}[section]
\newtheorem{proposition}[theorem]{Proposition}
\newtheorem{lemma}[theorem]{Lemma}
\theoremstyle{remark}\newtheorem{remark}[theorem]{Remark}
\DeclareMathOperator{\sgn}{sgn}
\DeclareMathOperator{\clip}{clip}
\newcommand{\R}{\mathbb R}
\newcommand{\ind}{\mathbf 1}
\newcommand{\Rbar}{\overline R}
\numberwithin{equation}{section}

\begin{document}
\begin{center}
{\color{ink}\LARGE\bfseries SRG Geometry\par}
\vspace{5pt}
{\Large Local Response, Neighbor Exchange,\par}
{\Large and Compression\par}
\vspace{10pt}
{\large Hilmir Fr\'imann Halld\'orsson\par}
\vspace{4pt}
{\small v0.1 \enspace$\cdot$\enspace 12 October 2026\par}
{\small Manuscript for GPT review\par}
\end{center}

\begin{abstract}
Moving points can change both the strength and the identity of their selected
neighbors. We analyze these two effects on the next amplitudes of the existing
Tower II model at a single fixed pass. A three-particle example gives an exact
response while the selected neighbors remain unchanged. We derive the
normalized-weight differential and a sufficient geometric margin preserving
selection. A seven-particle example isolates one fifth/sixth-neighbor exchange,
and a general jump formula retains the change in the selected enhanced-addition
term. The existing signed compressor preserves differences on a common
unclipped branch and identifies distinct nonzero limits exactly on common
signed plateaus. A fixed rescaling of the seven incoming amplitudes makes
the entire amplitude output constant despite a nonzero precompression jump.
This is not equality of complete moving-particle states. January's accepted
fixed-grid capture cells provide a separate comparison. All claims are
one-pass, conditional, exact-real deductions from unchanged source expressions;
no trajectories, new simulations or claims of physical validation are supplied.
\end{abstract}

\section{What changes when the points move?}\label{sec:question}

Imagine that each point receives values from its nearest neighbors.
Moving the points slightly can change the strength of those influences even
when the same neighbors remain selected. A further small movement can exchange
one selected neighbor for another. If the two candidates carry different
values, an arbitrarily small geometric change can then create a finite
amplitude difference. A final compression may preserve that difference or
map the two results to the same value.

We compare \emph{incoming configurations at the same incremented pass number}.
The configurations are not successive states produced by the model.
Their incoming amplitudes, frequencies and other parameters are held fixed
while geometry is compared, except in the one explicitly rescaled input
family of Section~\ref{sec:erasure}. Geometry comparison is not physical time.

\begin{center}
\begin{tabularx}{\linewidth}{@{}p{24mm}YY@{}}
\toprule
Model & Carrier and interaction & Meaning of equal output\\
\midrule
January & A finite periodic field; fixed index neighbors; selected signed mode.
& Equality of the complete field carrier.\\
Tower II & Moving particles with positions, amplitudes and velocities;
directed distance-selected neighbors.
& Equality of amplitudes alone leaves geometric state information separate.\\
\bottomrule
\end{tabularx}
\end{center}

These models are not composed into a replacement evolution. This paper presents
the accepted OC03 and OC04 deductions in a single exposition
\cite{OC03,OC04}. Their source is the frozen \emph{Golden Frequency Tower---Part II}
program \cite{T}; the January comparison reuses accepted OC02 results \cite{G}.
The historical reconstruction in SRG-I and SRG Origins remains distinct
\cite{SRGI,Origins}. The separate proposed SRG-II mechanism paper is not
renamed or replaced by this chapter.

\section{Three points: response with the same neighbors}\label{sec:three}

Place three particles on a line, with zero second and third coordinates:

\begin{center}
\begin{tabular}{@{}lccc@{}}
\toprule
Particle & Left & Middle & Right\\
\midrule
Position & $-1$ & $s$, \quad $-1<s<1$ & $+1$\\
Incoming amplitude & $1/10$ & $1/5$ & $3/10$\\
\bottomrule
\end{tabular}
\end{center}

There is one frequency band, with frequency $\omega=0.244$, no horizon and
no breathing of the compression bound. All other parameters keep their
defaults. The actual selector requests five neighbors; here it selects
\emph{both available other particles} in every row. It is not replaced by
a two-neighbor rule. These are specified input data, not the output of an
initializer \cite[lines 34, 62--103, 111--127]{T}.

At $s=0$ the middle particle receives equal influence from the two sides.
At $s=1/4$ the right particle is closer, so its larger amplitude receives
more weight. With the source regularizer $\delta=10^{-8}$, the two raw
influences on the middle particle are
\[
u_L=\frac{1}{1+s+\delta},\qquad
u_R=\frac{1}{1-s+\delta}.
\]
Their normalized right share and the resulting average are
\begin{align}
w_R&=\frac{u_R}{u_L+u_R}
=\frac{1+s+\delta}{2(1+\delta)}
=\frac12+\frac{s}{2(1+\delta)},\label{eq:three-weight}\\
\Rbar_M&=(1-w_R)\frac1{10}+w_R\frac3{10}
=\frac15+\frac{s}{10(1+\delta)}.\label{eq:three-mean}
\end{align}
Normalization changes the whole denominator; it cannot be held fixed while
the distances change.

Every off-diagonal distance is positive and at most two. Thus each row's
raw weight sum is at least $2/(2+\delta)>10^{-10}$ and the amplitude update
is active. Equal frequencies make the frequency factors one and fail the
strict enhanced-mixing guard $|\omega_i-\omega_j|>0.01$.
The enhanced coupling remains at its default value; the guard makes its
additions zero \cite[lines 156--188]{T}.

For the moment collect the sine and later H multipliers into $A_t$:
\begin{align}
H(t)&=2e^{-t/2}\cos t+e^{-3t/10}\sin t
       +\tfrac32e^{-7t/10}+e^{-2t/5},\label{eq:H}\\
a(t)&=1+0.01\sin(2\pi\omega t),\qquad
b_t=1+0.05H(t)/H_{\max},\qquad A_t=a(t)b_t.\label{eq:gains}
\end{align}
Here $t\ge1$ is the incremented integer pass. The retained exact bound
$|H(t)|\le9/2$ for $t\ge0$, with any fixed finite $H_{\max}\ge9/2$, gives
\begin{equation}
0.99\le a(t)\le1.01,\quad 0.95\le b_t\le1.05,\quad
0.9405\le A_t\le1.0605.\label{eq:gain-bounds}
\end{equation}
The source uses a sampled normalizer whose grid includes zero and then adds
$10^{-12}$. Its retained exact-real interpretation is $9/2+10^{-12}$.
We have not executed the normalizer or freshly certified its binary arithmetic.
Appendix~\ref{app:H} recalls the accepted bound without sampling.

Let $q$ mean the amplitude immediately before compression. With
$\beta=0.15$, the source's own-amplitude/neighbor combination gives
\begin{equation}
q_M(s)=A_t\left[(1-\beta)\frac15+\beta\Rbar_M(s)\right]
=A_t\left[\frac15+\frac{0.015s}{1+\delta}\right].
\label{eq:three-q}
\end{equation}
The full update order, and why positional stages leave this readout
unchanged, are given in Section~\ref{sec:contract}.

\begin{lemma}[Uniform branch bound]\label{lem:convex}
In this equal-frequency example, every precompression amplitude belongs to
$[0.09405,0.31815]$, throughout $-1<s<1$. Every output lies on the positive
unclipped branch of the source compressor.
\end{lemma}
\begin{proof}
Each normalized received average is a convex combination of amplitudes
between $0.1$ and $0.3$. Since $0<\beta<1$, combining it with the own
amplitude preserves this interval. Multiplying by
\eqref{eq:gain-bounds} gives
\begin{equation}
\frac{1881}{20000}=0.09405\le q_j
\le0.31815=\frac{6363}{20000}.\label{eq:uniform-q}
\end{equation}
The compressor has $L=9$ and $\varepsilon_T=10^{-9}$. Since $e>2$,
\[
e^{-9}<\frac1{512}<0.09405,\qquad
0.31815+\varepsilon_T<512<e^9.
\]
Therefore its logarithmic clip is inactive and its output is
$q_j+\varepsilon_T$. This includes epsilon in the branch check.
\end{proof}

Subtracting the two middle outputs cancels the same epsilon:
\begin{equation}
\boxed{R_M^{\rm next}(s)-R_M^{\rm next}(0)
       =\frac{0.015A_t}{1+\delta}s.}\label{eq:three-output}
\end{equation}
The coefficient is positive. This is a geometric response with no selected
neighbor exchange, established on the whole interval.

\begin{center}
\begin{tabular}{@{}rcc@{}}
\toprule
$s$ & $q_M(s)$ & $R_M^{\rm next}(s)-R_M^{\rm next}(0)$\\
\midrule
$-1/4$ & $0.204144213419495$ & $-0.003900844802657$\\
$0$ & $0.208045058222152$ & $0$\\
$1/4$ & $0.211945903024809$ & $+0.003900844802657$\\
\bottomrule
\end{tabular}
\end{center}

These rounded illustrations are transcribed from the saved OC03 scalar
results at the single common clock $t=1$, using $H_{\max}=9/2+10^{-12}$.
They were evaluated there at 80 decimal digits, without interval arithmetic,
and were not recomputed for this manuscript. The proof does not depend on
the number of saved examples.

\section{The source contract}\label{sec:contract}

For the local question, write the incoming positions as $x_i\in\R^3$ and
amplitudes as $R_i\in\R$. Assigned frequencies $\omega_i>0$, band identities,
the tower-frequency list and parameters are fixed. Positions and velocities
also belong to the particle state; the present readout is only its next
amplitude vector.

Without a horizon and without collisions, the source sorts distances in
each row, drops the zero self-entry and retains up to five other particles.
Selection is directed, not symmetrized. With a horizon $r_H$, eligible
particles satisfy the strict test $d_{ij}<r_H$. We write the selected
set as $\mathcal N_i$ and $d_{ij}=\|x_j-x_i\|$.

The pass order matters \cite[lines 129--215]{T}:
\begin{enumerate}
\item Increment the pass number and, if enabled, set the breathing bound from
incoming amplitudes. Breathing is false throughout the present analysis.
\item Store the incoming adjacency and distances, then move positions using
the selected neighbors.
\item Couple amplitudes using those stored distances and incidences.
Apply the sine multiplier, then the enhanced additive mixing.
\item Multiply amplitudes by the H factor.
\item Fuse positions within bands using the precompression amplitudes.
\item Compress amplitudes, then add positional noise and record velocities.
\end{enumerate}
In particular, within-pass spatial movement does not refresh the weights
used by that same amplitude calculation. Fusion changes positions only, and
noise is added only to positions. Neither modifies the amplitude readout
after its indicated stage.

For selected neighbors, the actual raw weights and normalization are
\begin{align}
\chi_{ij}&=\exp[-c(\log(\omega_j/\omega_i))^2],\qquad
u_{ij}=\frac{\chi_{ij}}{d_{ij}+\delta},\label{eq:raw-weights}\\
S_i&=\sum_{j\in\mathcal N_i}u_{ij},\qquad
w_{ij}=\frac{u_{ij}}{S_i},\qquad
\Rbar_i=\sum_{j\in\mathcal N_i}w_{ij}R_j.\label{eq:normalized}
\end{align}
The defaults are $c=3$, $\delta=10^{-8}$, $\beta=0.15$,
$L=9$ and $\varepsilon_T=10^{-9}$.
A row with no neighbors or with $S_i<10^{-10}$ is skipped at the mixing
stage: its sine and enhanced additions are both skipped.
Equality $S_i=10^{-10}$ is active in the code, but provides no open active
margin. Our local theorems require strict margins above that threshold
\cite[lines 151--163]{T}; see also \cite[A03]{A}.

For an active row define $a_i(t)=1+0.01\sin(2\pi\omega_i t)$ and use $b_t$
from \eqref{eq:gains}. The source order gives
\begin{equation}
q_i=b_t\left\{
a_i(t)\big[(1-\beta)R_i+\beta\Rbar_i\big]
+\sum_{j\in\mathcal N_i}\eta_{ij}R_j
\right\}.\label{eq:active}
\end{equation}
The coefficients $\eta_{ij}$ are determined by the source frequency tests,
not fitted coefficients. Appendix~\ref{app:eta} gives their full expression,
including every strict guard and loop multiplicity.
The sine does \emph{not} multiply the later additive sum.
For a skipped row, the corresponding precompression amplitude is instead
$b_tR_i$; we do not extend \eqref{eq:active} through such a boundary.

\section{Smooth response inside a neighbor region}\label{sec:smooth}

Let positions depend differentiably on a geometry parameter $\tau$.
A dot denotes differentiation with respect to $\tau$, not time.
Hold all incoming amplitudes, assigned frequencies, bands, tower list,
clock and parameters fixed. Work on an open region with nonzero pair
distances, fixed selected incidences and $S_i>10^{-10}$.

\begin{theorem}[Normalized geometry response, OC03]\label{thm:smooth}
For any active finite selected row,
\begin{align}
\dot{\Rbar}_i
&=-\sum_{j\in\mathcal N_i}w_{ij}(R_j-\Rbar_i)
          \frac{\dot d_{ij}}{d_{ij}+\delta},\label{eq:derivative}\\
\dot d_{ij}
&=\frac{x_j-x_i}{d_{ij}}\cdot(\dot x_j-\dot x_i),\label{eq:distance-diff}\\
\dot q_i&=b_ta_i(t)\beta\,\dot{\Rbar}_i.\label{eq:q-diff}
\end{align}
These are statements about the conditional precompression map. The final
amplitude derivative additionally requires a smooth compressor branch.
\end{theorem}
\begin{proof}
Put $\ell_{ij}=\dot d_{ij}/(d_{ij}+\delta)$. The fixed frequencies give
$\dot u_{ij}=-u_{ij}\ell_{ij}$ and
$\dot S_i=-S_i\sum_k w_{ik}\ell_{ik}$. Differentiating the shared denominator,
\begin{equation}
\dot w_{ij}
=\frac{\dot u_{ij}S_i-u_{ij}\dot S_i}{S_i^2}
=-w_{ij}\left(\ell_{ij}-\sum_k w_{ik}\ell_{ik}\right).
\label{eq:weight-diff}
\end{equation}
Multiplying by $R_j$ and summing yields \eqref{eq:derivative}, since the
amplitudes are fixed. Differentiating
$d_{ij}^2=(x_j-x_i)\cdot(x_j-x_i)$ gives \eqref{eq:distance-diff}.
All $\eta_{ij}$ and their selected incidences are fixed in this region.
Differentiation of \eqref{eq:active} therefore gives \eqref{eq:q-diff}.
\end{proof}

If every selected neighbor has the same amplitude, every contrast
$R_j-\Rbar_i$ vanishes. Unequal amplitudes permit a response but do not
guarantee one. For example, when all selected $\ell_{ij}$ equal one common
value, the response vanishes because
$\sum_jw_{ij}(R_j-\Rbar_i)=0$. Common translation is a special case with
all distance changes zero.

\begin{proposition}[A sufficient selection margin, OC03]\label{prop:margin}
Suppose a row has at least six other eligible candidates, fixed horizon
membership and fifth/sixth distance gap
$\Delta_i=d_{i,(6)}-d_{i,(5)}>0$.
If each particle displacement has norm at most $\rho$, then
\begin{equation}
4\rho<\Delta_i\label{eq:margin}
\end{equation}
preserves its selected set, provided collisions and horizon/active-guard
crossings are independently excluded.
\end{proposition}
\begin{proof}
For perturbations $e_i$, the reverse triangle inequality gives
\[
|d'_{ij}-d_{ij}|\le\|e_j-e_i\|\le2\rho.
\]
For any originally selected $j$ and excluded eligible $k$,
\[
d'_{ik}-d'_{ij}
\ge(d_{i,(6)}-2\rho)-(d_{i,(5)}+2\rho)
=\Delta_i-4\rho>0.
\]
Thus no excluded candidate can overtake a selected one.
\end{proof}

Useful independent sufficient margins are
\begin{align}
2\rho&<\min_{i\ne j}d_{ij}
&&\text{(no collisions)},\label{eq:collision}\\
2\rho&<\min_{i\ne j}|d_{ij}-r_H|
&&\text{(fixed strict horizon membership)},\label{eq:horizon}\\
\sum_{j\in\mathcal N_i}
\frac{\chi_{ij}}{d_{ij}+\delta+2\rho}&>10^{-10}
&&\text{(active row)}.\label{eq:guard}
\end{align}
The horizon condition is needed only when a horizon exists.
The last expression is a lower bound on the perturbed raw weight sum
once incidences are preserved. A geometric cutoff gap alone does not
ensure activity.
The condition $4\rho<\Delta_i$ is sufficient, not necessary or claimed sharp.
With at most five eligible other particles, all are selected; no sixth
distance exists and no cutoff gap is needed. Ties within a selected set
are harmless to its incidence set, unlike a tie across its cutoff.

\section{Seven points: an isolated neighbor exchange}\label{sec:seven}

Keep a focal particle at zero and four common neighbors to its right at
$1/2,1,3/2,2$. Two farther particles compete for the last selected place.
The right candidate A, near $+3$, has amplitude $0.1$; the left candidate B,
near $-3$, has amplitude $0.3$. Move both candidates slightly left: A is
closer to the focal point. Move both slightly right: B is closer.
Equal distances at the boundary do not make their incoming amplitudes equal.

\begin{center}
\begin{tabular}{@{}lccccccc@{}}
\toprule
Particle & B & i & c1 & c2 & c3 & c4 & A\\
\midrule
Position & $-3+s$ & $0$ & $1/2$ & $1$ & $3/2$ & $2$ & $3+s$\\
Amplitude & $3/10$ & $1/5$ & $1/5$ & $1/5$ & $1/5$ & $1/5$ & $1/10$\\
\bottomrule
\end{tabular}
\end{center}

All second and third coordinates are zero, $-1/4<s<1/4$, and the seven
particles have one band at frequency $0.244$, no horizon, no breathing,
and unchanged default parameters. This is a specified incoming family.

\begin{proposition}[Isolation in the full directed graph, OC04]\label{prop:isolation}
For $s<0$ the focal row selects A and excludes B; for $s>0$ it selects B
and excludes A. At zero these candidates tie for fifth/sixth place.
Every other row retains its selected set throughout the interval.
Every row has a strict active-weight margin.
\end{proposition}
\begin{proof}
The focal common distances are $1/2,1,3/2,2$. The candidate distances
$3+s$ and $3-s$ are both greater than $11/4>2$; their difference is $2s$,
with slope 2. They exchange order transversely without a collision.

At any common coordinate $x\in\{1/2,1,3/2,2\}$, B is farther than every
focal/common particle, since $x+3-s>13/4$ while those distances are at most
2. B is also farther than A, because
\[
(x+3-s)-(3+s-x)=2(x-s)>0.
\]
Hence every common row uniquely excludes B. The A--B distance is 6.
All other distances from A are at most $3+s<13/4$, so A excludes B;
all other distances from B are at most $5-s<21/4$, so B excludes A.
All positions are distinct and all pair distances are at most 6.
Each row selects five equal-frequency inverse-distance weights, giving
\begin{equation}
S_i\ge\frac5{6+\delta}
=\frac{500000000}{600000001}>10^{-10}.\label{eq:seven-guard}
\end{equation}
This proves isolation of the focal exchange and strict activity on the
whole interval.
\end{proof}

At the switch define
\begin{equation}
S=\sum_{k=1}^4\frac1{k/2+\delta},\qquad
u_*=\frac1{3+\delta},\qquad w_*=\frac{u_*}{S+u_*}.
\label{eq:fifth-share}
\end{equation}
The last quantity is the normalized share of the fifth neighbor. Equal
frequencies again turn off the guarded enhanced additions without changing
their parameters. The limiting received averages are
$1/5-w_*/10$ and $1/5+w_*/10$. With $A_t$ from \eqref{eq:gains},
\begin{align}
q_-&=A_t\left(\frac15-\frac{3}{200}w_*\right),&
q_+&=A_t\left(\frac15+\frac{3}{200}w_*\right),\label{eq:focal-limits}\\
\Delta q_i&=q_+-q_-=\frac3{100}A_tw_*>0.\label{eq:focal-jump}
\end{align}
All seven incoming amplitudes lie in $[0.1,0.3]$, so the convex argument
of Lemma~\ref{lem:convex} applies unchanged. All outputs are positive and
unclipped. The same epsilon cancels and the compressed jump equals
\eqref{eq:focal-jump}.

The regularized share is approximately $0.0740740747828074933$.
Only the formal $\delta\to0$ comparison gives $S=25/6$, $u_*=1/3$ and
$w_*=2/27$. We do not substitute that comparison for the actual source weight.
The retained $t=1$ scalar results give
\[
q_-\approx0.206889252332082,\qquad
q_+\approx0.209200864112223,\qquad
\Delta q_i\approx0.002311611780141188.
\]

The source specifies a sorting call, not a separately defined portable
exact-tie rule. No tie-sorting experiment is needed or supplied.
The two one-sided limits are independent of the tie choice. Because they
differ, no choice of the single tie value can make this amplitude map
continuous there.

\subsection{The general exchange law}\label{sec:exchange}

Let the focal selected sets on two sides of a boundary be
$C\cup\{A\}$ and $C\cup\{B\}$, where $C$ contains four unchanged neighbors.
Assume $d_{iA}=d_{iB}=d_*>0$ at the boundary and common limits for other
distances. Keep the frequencies, incoming amplitudes, bands, tower list,
parameters and clock fixed. Assume a transverse candidate-order exchange,
no further simultaneous focal cutoff, horizon or guard event, and both
limiting raw weight sums strictly above $10^{-10}$.
These margins allow \eqref{eq:active} on both sides.

At the boundary put $u_j=\chi_{ij}/(d_{ij}+\delta)$,
$S=\sum_{j\in C}u_j>0$ and $V=\sum_{j\in C}u_jR_j$.
Normalize separately:
\begin{equation}
\Rbar_-=\frac{V+u_AR_A}{S+u_A},\qquad
\Rbar_+=\frac{V+u_BR_B}{S+u_B}.\label{eq:exchange-means}
\end{equation}

\begin{theorem}[Active-row exchange, OC04]\label{thm:exchange}
Under these assumptions the one-sided differences satisfy
\begin{align}
\Delta\Rbar_i
&=\frac{u_B}{S+u_B}\left(R_B-\frac VS\right)
-\frac{u_A}{S+u_A}\left(R_A-\frac VS\right),\label{eq:contrast}\\
\Delta q_i
&=b_t\left[a_i(t)\beta\,\Delta\Rbar_i
           +\eta_{iB}R_B-\eta_{iA}R_A\right].\label{eq:exchange-jump}
\end{align}
Since $b_t>0$ under the retained defaults, the exact zero-jump condition is
\begin{equation}
a_i(t)\beta\left[
\frac{V+u_BR_B}{S+u_B}-\frac{V+u_AR_A}{S+u_A}
\right]+\eta_{iB}R_B-\eta_{iA}R_A=0.\label{eq:zero-jump}
\end{equation}
\end{theorem}
\begin{proof}
Subtract $V/S$ from each fraction in \eqref{eq:exchange-means}; this gives
\eqref{eq:contrast}. Subtract the two source-ordered expressions
\eqref{eq:active}. The own-amplitude baseline and every common-neighbor
additive term cancel, leaving \eqref{eq:exchange-jump}. The lower bound
$b_t\ge0.95$ then gives \eqref{eq:zero-jump}.
\end{proof}

An additive coefficient is independent of distance but its \emph{selected
identity} changes here. Its zero positional derivative inside a fixed
neighbor region does not justify discarding the selected term at a switch.
This is the boundary distinction between Theorems~\ref{thm:smooth}
and~\ref{thm:exchange}.

Equal candidate distances need not give $u_A=u_B$, because the frequency
matches may differ. Thus the two denominators generally differ.
If $R_A=R_B=r$, the received-average contrast is
\[
\left(\frac{u_B}{S+u_B}-\frac{u_A}{S+u_A}\right)
\left(r-\frac VS\right),
\]
which need not vanish. Even when all selected incoming amplitudes equal
$r>0$, the remaining jump can be
$b_tr(\eta_{iB}-\eta_{iA})$. A same-frequency A has coefficient zero,
while an available adjacent $\varphi$-band B has a strictly positive
coefficient under the source guard.

Conversely, unequal amplitudes do not guarantee a jump. For example,
when $V=0$, define
$K_A=a_i\beta u_A/(S+u_A)+\eta_{iA}>0$ and similarly $K_B>0$.
Whenever $K_A\ne K_B$, choosing $R_B=K_AR_A/K_B$ with $R_A>0$ makes
\eqref{eq:zero-jump} zero despite unequal candidates. This is a conditional
algebraic illustration, not an extra numerical fixture.
In the equal-frequency example, $\eta=0$ and $u_A=u_B$, so the law reduces to
$\Delta q_i=A_t\beta w_*(R_B-R_A)$ and gives \eqref{eq:focal-jump}.

\section{What the existing compressor identifies}\label{sec:compression}

The frozen source applies
\begin{equation}
C_T(q)=\sgn(q)\exp\!\left[
\clip\big(\log(|q|+\varepsilon_T),-9,9\big)\right],
\qquad \varepsilon_T=10^{-9}.\label{eq:compressor}
\end{equation}
Set $u_T=e^{-9}$, $v_T=e^9$, $f_T=u_T-\varepsilon_T>0$ and
$h_T=v_T-\varepsilon_T$. Monotonicity of logarithm and exponential gives
the complete exact-real branches:
\begin{equation}
C_T(q)=
\begin{cases}
0,&q=0,\\
\sgn(q)u_T,&0<|q|\le f_T,\\
q+\varepsilon_T\sgn(q),&f_T<|q|<h_T,\\
\sgn(q)v_T,&|q|\ge h_T.
\end{cases}\label{eq:branches}
\end{equation}
The floor upper edge and ceiling lower edge are included; zero is a
separate case \cite[lines 42--44, 208--209]{T}.

\begin{theorem}[Nonzero limiting inputs, OC04]\label{thm:compression}
For finite nonzero one-sided precompression limits $q_-,q_+$, the compressed
jump is $C_T(q_+)-C_T(q_-)$. A common signed unclipped branch preserves the
input difference exactly. Distinct such limits have equal outputs if and
only if both lie on one common signed floor or ceiling plateau.
Opposite nonzero signs cannot have the same output.
\end{theorem}
\begin{proof}
The branches are continuous at every nonzero input. At the positive floor
edge the middle expression gives $f_T+\varepsilon_T=u_T$; at the positive
ceiling edge it gives $h_T+\varepsilon_T=v_T$. Odd reflection gives the
negative edges. Hence continuity passes each nonzero limit through $C_T$.
The clipping edges are nondifferentiable, not discontinuous.

On either fixed signed unclipped branch the same signed epsilon cancels
under subtraction. On the positive side, the fibers of $u_T$ and $v_T$
are exactly $(0,f_T]$ and $[h_T,\infty)$; between them the function is
strictly increasing and has singleton fibers. The negative plateau fibers
are $[-f_T,0)$ and $(-\infty,-h_T]$, with singleton fibers in between.
The output signs separate the two sides. These exhaust the possibilities.
\end{proof}

For smooth incoming comparisons, the final amplitude derivative is
$\dot q_i$ on a nonzero unclipped branch and zero on the interior of a
fixed signed plateau. This statement does not extend a derivative through
the clipping edges or zero.

\subsection{Zero requires an approach condition}

If precompression values tend to zero and are eventually strictly positive,
they eventually belong to $(0,f_T]$, so their compressed limit is $+u_T$.
Eventually strictly negative approach gives $-u_T$. Values identically
zero on a side near the boundary remain zero.
Thus $C_T(0)=0$ cannot be substituted for an unqualified limit.
If none of these approach conditions is established, a zero precompression
limit alone does not establish a compressed limit. We make no general
oscillatory-boundary classification.

\subsection{A whole amplitude vector on one floor}\label{sec:erasure}

The unscaled seven-point family has the positive jump
\eqref{eq:focal-jump}, preserved by its unclipped branch. Now keep the same
positions, selected rule, frequency, clock and model parameters but multiply
every incoming amplitude by the single fixed data scale $\kappa=10^{-4}$.
This is a second specified input family, not a parameter sweep or a
modified operator.

The equal-frequency additions still vanish and the remaining amplitude
expression is linear in incoming amplitudes. Because $L$ does not breathe,
precompression amplitudes scale by $\kappa$ without changing the compressor.
The uniform bound \eqref{eq:uniform-q} therefore gives
\begin{equation}
0<q_j^{(\kappa)}(s)\le\kappa\,\frac{6363}{20000},\qquad
\kappa\,\frac{6363}{20000}+10^{-9}
=\frac{3977}{125000000}<\frac1{19683}<e^{-9}.
\label{eq:scaled-bound}
\end{equation}
The first strict difference is the positive rational number
$46720709/2460375000000$; the last inequality follows from $e<3$.
For completeness, $e=\sum_{n\ge0}1/n!>2$, and
$n!\ge2^{n-1}$ for $n\ge2$, strictly for $n\ge3$, gives
$e<2+\sum_{n\ge2}2^{-(n-1)}=3$.

Every row is therefore strictly on the positive floor throughout the
whole interval, including either focal choice at the tie:
\begin{equation}
\boxed{(R_1^{\rm next},\ldots,R_7^{\rm next})
       =(e^{-9},\ldots,e^{-9}).}\label{eq:whole-vector}
\end{equation}
The focal precompression jump remains the strictly positive
$\kappa(3/100)A_tw_*$, but the compressed jump is zero.
The saved $t=1$ values are
$q_-^{(\kappa)}\approx0.0000206889252332082$ and
$q_+^{(\kappa)}\approx0.0000209200864112223$; their difference is approximately
$2.311611780141188\times10^{-7}$ before compression and zero afterward.
These retained values illustrate the uniform proof; they are not a
seven-particle update.

Equation~\eqref{eq:whole-vector} is \emph{not full particle-state capture}.
Band fusion has already used the precompression amplitudes through both
its weighted center and its positional movement factor.
Positions and velocities carry separate information.
No equality or inequality of complete final particle states, and no later
trajectory consequence, follows automatically from this amplitude theorem.

\section{January: a separate fixed-grid reference}\label{sec:january}

January's accepted signed-mode analysis has a finite periodic field
$X\in\R^{NM}$ as its carrier. This signed mode is selected for the analysis;
the source app initially selects paper mode. Its fixed index stencil is
\begin{align}
(\mathcal LX)_{ij}
&=X_{i+1,j}-2X_{ij}+X_{i-1,j}\notag\\
&\quad+\frac{25}{9}(X_{i,j+1}-2X_{ij}+X_{i,j-1}),\label{eq:grid-stencil}\\
T_G&=I+\frac1{20}\mathcal L,\qquad
P(X)=T_GX+\frac1{10}\tanh X,\qquad F_m(X)=C_s(mP(X)).
\label{eq:grid-P}
\end{align}
Here $T_G$ denotes the operator called $T$ in OC02; the subscript avoids
confusion with the Tower II source key. The tanh and compressor act sitewise.
Duplicate incidences for axes of length two are counted twice; on an axis
of length one the corresponding neighbor is the site itself. This is the
original periodic index stencil, not an embedded-surface Laplacian.
The constants are $L=10$ and $\varepsilon_G=10^{-6}$:
\[
C_s(q)=\sgn(q)\clip(|q|+\varepsilon_G,e^{-10},e^{10})
\quad(q\ne0),\qquad C_s(0)=0.
\]
These definitions are retained from OC02 \cite[equations (1)--(2)]{G}
and its source correspondence \cite{January}.

For a sign field $\sigma_{ij}\in\{-1,1\}$ and a fixed gain $m>0$, the
accepted one-step capture condition is
\begin{equation}
F_m(X)=e^{-10}\sigma
\quad\Longleftrightarrow\quad
0<\sigma_{ij}mP(X)_{ij}\le e^{-10}-10^{-6}
\quad\text{for every site}.\label{eq:grid-capture}
\end{equation}
Indeed, the signed floor has exactly that nonzero half-open scalar
preimage at each site. The zero face is excluded; the upper face is included.

OC02's retained Theorem 2 states that $P$ is a global smooth diffeomorphism.
Its mechanism is that the symmetric stencil has lower quadratic-form bound
$11/45$, so $DP=T_G+\operatorname{diag}(\operatorname{sech}^2X)/10$ is
uniformly positive definite. The strictly convex, coercive potential
$X^TT_GX/2+\sum_{ij}\log\cosh(X_{ij})/10$ supplies a unique solution of
$P(X)=z$ for every $z$, and the inverse function theorem supplies the smooth
inverse. This recalls the accepted argument, not a new classification.

Writing $b=(e^{-10}-10^{-6})/m$, the capture set is consequently
\[
P^{-1}\{z:0<\sigma_{ij}z_{ij}\le b\text{ at every site}\},
\]
the inverse image of a half-open rectangular cell. The interior uses strict
inequalities on both faces; the closure includes zero.
A field may reach an arbitrary signed floor target in one step without that
target persisting. Persistence under subsequent gains requires the accepted
common-fixed-pattern conditions. We do not recount that classification or
compute multistep preimages \cite[lines 224--269]{G}.

Both models can erase a distinction at compression. Equality of January's
complete field carrier and equality of Tower II's amplitude vector are,
however, different mathematical statements. The latter leaves moving
geometry outside the equality. No January basin or fixed-pattern theorem
has been transferred to particle states.

\section{Scope, evidence and interpretation}\label{sec:limits}

The results concern exact-real algebra of the unchanged source expressions
at one fixed pass. Smooth response is conditional on an active region with
unchanged incidences and nonzero distances; the exchange theorem concerns
one isolated focal event with strict limiting margins. Final readout
statements retain the compressor's sign and branch conditions.
Not every exchange jumps, and compression does not erase every difference.

The H bound and January results are inherited mathematics. The geometry
response, selection margin, exchange law and constructive compressor examples
are OC03/OC04 deductions in this research program. Their proofs are reproduced
here. No literature survey has established global novelty or priority over
general results on normalized weights, nearest-neighbor maps or clipping.
No physical validation is claimed.

OC03 and OC04 each retain eight mathematical or scalar groups and two
groups for input identity and preservation. Their saved exact rational and symbolic
identities and 80-digit, non-interval scalar corroborations support examples;
they do not replace the proofs. Neither checker was rerun for this paper.
GPT's separate OC04 review reports independent symbolic and 100-digit
scalar corroboration, without executing either checker or historical program.
That review is not external peer review, a binary execution certificate or
a replacement for the saved evidence.

No source normalizer, initializer, particle update, trajectory or stochastic
step was executed in preparing this edition. New QA concerns transcription,
source mapping, typography and the review artifacts only.
There are no arrival statistics, autonomous-attractor results, general
recursion theorem or full-state capture theorem here.
Later-pass questions remain future work; they do not weaken the stated
one-pass conclusions.

\subsection*{Conclusion}

Incoming geometry affects the next amplitude in two distinct ways:
smooth variation of normalized influences within a selected-neighbor region,
and a possible finite jump when the selected identity changes.
The jump law keeps both the normalized average and the switched enhanced
addition. The existing compressor then preserves a common signed unclipped
difference or identifies a common signed plateau.
The two explicit geometric examples make these alternatives checkable.
They establish a local amplitude account while keeping the complete
moving-particle state, later dynamics and January's separate field carrier
in their proper scopes.

\subsection*{Assistance, status and reuse}

Hilmir Fr\'imann Halld\'orsson is the named human author.
GPT supplied mathematical review of OC03/OC04 and the manuscript work order.
Codex performed the recorded OC03/OC04 analysis and finite checks, then
assembled this manuscript, checked its source correspondence and completed
artifact QA. Those assistance roles are not coauthorship or external peer
review. The underlying analyses are internally accepted; this new v0.1
manuscript and its proposed public payload await GPT review.
No institutional affiliation, DOI or series number is assigned.

This edition grants no new blanket reuse licence. The software repository's
Apache licence does not automatically cover scientific research.
Linked deposited sources retain their source-specific notices, including
the identified CC BY 4.0 notices. Those notices are not reassigned to the
derived manuscript or its reconstruction records.

\appendix
\section{The complete enhanced-addition coefficient}\label{app:eta}

Let $\nu=0.001$, $\varphi=(1+\sqrt5)/2$,
$D_{ij}=|\omega_i-\omega_j|$, $U_{ij}=\omega_i+\omega_j$, and let
$\mathcal W$ be the fixed source tower-frequency list.
The coefficient in \eqref{eq:active} is exactly
\begin{equation}
\begin{aligned}
\eta_{ij}={}&\nu\ind_{\{D_{ij}>0.01\}}
\Bigg[
\sum_{p\in\{-2,-1,1,2\}}
e^{-5[\log(\omega_i/\omega_j)-p\log\varphi]^2}\\
&\quad+\frac12\sum_{\Omega\in\mathcal W}
\left(
\ind_{\{|U_{ij}-\Omega|/\Omega<0.2\}}
+\ind_{\{D_{ij}>0.01\}}
 \ind_{\{|D_{ij}-\Omega|/\Omega<0.2\}}
\right)
\Bigg].
\end{aligned}\label{eq:eta}
\end{equation}
This transcribes lines 169--188 of the frozen source \cite{T}.
Each ratio test and each matching tower entry contributes separately.
The repeated difference guard records the inner test; it is redundant
once the outer guard holds. The whole coefficient is added after sine
multiplication, as in \eqref{eq:active}.
For identical frequencies the outer guard fails and $\eta_{ij}=0$.
Within fixed incidences the coefficient has no positional derivative,
but a changed selected identity contributes the last two terms of
\eqref{eq:exchange-jump}.

\section{The retained H bound}\label{app:H}

For the same expression \eqref{eq:H}, the accepted argument
\cite[lines 185--195]{G} uses $|\sin t|\le t$ to bound
\[
|H(t)|\le B(t)+\tfrac32e^{-7t/10}+e^{-2t/5},
\qquad B(t)=2e^{-t/2}+te^{-3t/10}.
\]
Its derivative is
$B'(t)=-e^{-t/2}+(1-3t/10)e^{-3t/10}$.
For $t\ge10/3$ both summands are nonpositive.
For $0\le t<10/3$,
$1-3t/10\le1-t/5\le e^{-t/5}$, so again $B'(t)\le0$.
Thus $B(t)\le2$ and $|H(t)|\le2+3/2+1=9/2$, attained at zero.
The source's 50,000-point grid on $[0,100]$ includes zero; its exact-real
maximum is therefore $9/2$, and its added guard gives $H_{\max}=9/2+10^{-12}$.
This inherited mathematical argument makes no claim of a fresh
floating-point audit of that helper.

\begin{thebibliography}{9}
\raggedright

\bibitem{T}
\emph{The Golden Frequency Tower---Part II}.
Historical source, 14 March 2026, file
\texttt{srg\_golden\_tower\_II.py}.
Header credit: Claude (collaborating with pzychozen).
Frozen 34,980-byte version; relevant physical lines 34, 37--46,
62--103, 111--127 and 129--215.
Recorded public route:
\href{https://zenodo.org/records/19023429/files/simulations.zip}{Zenodo 19023429, simulations archive}.
Source-specific CC BY 4.0 notice retained from the existing public source index.

\bibitem{OC03}
SRG research records.
\emph{OC03: Geometry Response with Unchanged Neighbors}.
12 October 2026. Accepted internal mathematical analysis,
Sections 1--5, equations B1--B5, C1--C6, D1--D3 and E1--E2.
The full internal report is not redistributed in this edition.

\bibitem{OC04}
SRG research records.
\emph{OC04: A Neighbor Exchange Can Survive or Disappear under Compression}.
12 October 2026. Accepted internal mathematical analysis,
Sections 1--5, equations (1)--(20).
GPT acceptance: 12 October 2026, within its one-pass exact-real scope.
The full internal report and review are not public payload files.

\bibitem{G}
SRG research records.
\emph{Signed Floor Patterns and Exact One Step Capture} (OC02).
11 October 2026. Accepted internal report.
Definitions at physical lines 27--63; H bound at 185--195;
Theorem 2 and capture cells at 224--269, equations (12)--(13).

\bibitem{A}
SRG research records.
\emph{Source and Notation Addendum}.
11 October 2026. A03, physical lines 33--47:
inactive-row branches and sine-before-additive order.
Internal retained source correspondence.

\bibitem{January}
Hilmir Fr\'imann Halld\'orsson.
\emph{An Experimental Framework for Exploring Recursive Numerical Operators:
TGMO, RPCO, and REFU as a Reproducible Computational Laboratory}.
25 January 2026, physical pp.\ 2--3; accompanying
\texttt{app.py}, \texttt{tgmo.py} and \texttt{refu.py}.
Recorded exact-version route:
\href{https://zenodo.org/records/18285147}{Zenodo 18285147}.
The paper's cover title differs from the deposit title.
Source-specific CC BY 4.0 notices remain with these works.

\bibitem{SRGI}
Hilmir Fr\'imann Halld\'orsson.
\emph{SRG-I: Foundational Reconstruction}, v0.1.1, 10 October 2026.
Historical reconstruction, recursive well-posedness and dual-sector symmetries.
\href{https://github.com/pzychozen/trioctagon-physics/tree/3e5972be2f233b6c8ec3c2ebe0687ee1f8f2ff1f/SRG_REDISCOVERY/PAPERS/SRG_I_FOUNDATIONAL_RECONSTRUCTION/publication/v0.1.1}{Repository edition at the supplied publication checkpoint}.

\bibitem{Origins}
Hilmir Fr\'imann Halld\'orsson.
\emph{SRG Origins---From Geometry to Recursive Operators}, v0.1,
11 October 2026.
\href{https://github.com/pzychozen/trioctagon-physics/tree/3e5972be2f233b6c8ec3c2ebe0687ee1f8f2ff1f/SRG_REDISCOVERY/PAPERS/SRG_ORIGINS_FROM_GEOMETRY_TO_RECURSION/publication/v0.1}{Repository edition at the supplied publication checkpoint}.
Source access, exact versions and source-specific rights are recorded
in its source index.

\end{thebibliography}

\small
Source routes above are carried from locally inspected public indexes,
not newly checked online. This edition reproduces the definitions and proofs
needed for its central claims. Its companion source index records exact
identities and the access limits of retained internal reports.
\end{document}
